\begin{definition}[Smoothness away from critical points, as a hypothesis on $p$]
Let $p\in\mathbb{R}$.  We say that \emph{$p$ satisfies smoothness away from critical points}
when the following holds.  Let $\Omega\subseteq\mathrm{Plane}$, let $u$ be $p$-harmonic in
$\Omega$ with weak gradient $g$ (which in particular makes $\Omega$ open), and let
$U\subseteq\Omega$ be open.  If $u$
is of class $C^1$ on $U$ and its classical gradient satisfies $Du(x)\neq0$ for every
$x\in U$, then $u$ is of class $C^\infty$ on $U$.

This is the statement the paper uses when it asserts, in the proof of its estimate at
noncritical points, that ``since $Dv$ does not vanish in $B_{2\tau}$, the function $v$ is
smooth there''; the reference given is Lindqvist, \emph{Notes on the $p$-Laplace equation},
Report 102, University of Jyv\"askyl\"a, 2006, Section~4.  The equation
$-\operatorname{div}(\lvert Du\rvert^{p-2}Du)=0$ is uniformly elliptic with smooth
coefficients on any region where $Du$ is bounded away from $0$ and from $\infty$, so
classical elliptic regularity applies there.

The hypothesis is stated as a property of $p$ rather than proved: it is one of the results
the paper cites from the literature.
\end{definition}
